\section{Related Works}

The current landscape of Text-to-Query research is marked by a variety of approaches that leverage different techniques and methodologies to address typical problems encountered in this class of tasks. Broadly speaking, we can categorize existing works into two main lines of research: (i) the development of benchmark datasets specifically designed for training and evaluating Text-to-Query models, and (ii) the design of general Text-to-Query solutions that can be adapted to various types of databases and query languages. For our purposes, in the first line of research, we are mainly interested in datasets focused on the Text-to-Cypher task, while in the second line of research, we will explore more general Text-to-Query solutions that can be adapted to our specific task.

\subsection{Text-to-Cypher Datasets} 

An example include the \href{https://huggingface.co/datasets/neo4j/text2cypher-2024v1}{neo4j/text2cypher-2024v1} dataset introduced by the Neo4j-Labs team \cite{neo4jlabs-text2cypher-2024}, which brings together instances from publicly available datasets, cleaning and organizing them for smoother use. This dataset consists of 44,387 instances, with 39,554 instances in the training split and 4,833 instances in the test split where for each instance, a natural language question is paired with: its corresponding Cypher query, the entire schema of the underlying graph database and, if available, an identifier to the specific Neo4j database hosting the entire knowledge graph from which the instance was derived. The Neo4j-Labs team, in fact, provide a remote access to multiple already deployed Neo4j graph databases, freely and publicly accessible, that can be used to execute the Cypher queries present in the dataset. This important aspect paired with the vast diversity of the included knowledge graphs (about 17 different Neo4j graph databases are used to derive the instances in the dataset), makes it well-suited for our purposes.

Another notable example is the \href{https://huggingface.co/datasets/megagonlabs/cypherbench}{megagonlabs/cypherbench} dataset introduced by Megagon Labs \cite{cypherbench-megagonlabs-2024} with 10,882 instances for training and evaluating text-to-Cypher models. To create such a benchmark, the authors first developed a novel RDF-to-Property-Graph transformation engine that converts RDF data into high quality locally deployable property knowledge graphs while preserving the original semantics. Then, they applied this engine to Wikidata, generating 11 large-scale multi-domain property graphs that captures a wide range of real-world entities and relationships. Lastly, they designed a synthetic data generation pipeline that, starting from graph patterns templates, automatically creates diverse natural language questions and their corresponding Cypher queries based on the structure and content of the generated property graph. The peculiarity of this approach lies in its focus on comprehensiveness, in graph patterns covered, and diversity, in graph schemas and use cases addressed, of the resulting benchmark dataset. However, the lack of already deployed, publicly accessible and remote Neo4j graph databases endpoints, force anyone willing to use this dataset to first locally deploy the generated property graphs before being able to execute the Cypher queries present in the dataset. To further complicate things, the configuration method of such local deployments provided by the authors requires significant hardware requirements (at least 64 GB of RAM when deploying the 7 test graphs and 128GB when deploying all 11 graphs\footnote{This was explicitely stated in the 3rd step of the quick start of the \texttt{README.md} file in their repo \url{https://github.com/megagonlabs/cypherbench}}) that are out of reach for many researchers. For this reason, we decided to not use this dataset for our experiments.

\subsection{General Text-to-Query Solutions}
\label{sec:text-to-query-solutions}

Several solutions have been proposed to tackle the general Text-to-Query translation task, primarily leveraging transformer-based architectures \cite{vaswani-etal-2017-attention}, both decoder-only and encoder-decoder models, in combination with various techniques. %
%
To address problem \ref{itm:representing} in Section \ref{sec:introduction} and ensure the generation of \textbf{schema-compliant}, syntactically and semantically correct queries, researchers have extensively adopted \textbf{in-context learning} (ICL) combined with \textbf{prompt engineering}. These techniques exploit the few-shot and zero-shot capabilities of large language models by providing carefully crafted prompts containing \textit{task instructions}, \textit{few-shot demonstrations} (retrieved from the training set), and relevant contextual information such as \textit{schema elements} to guide query generation \cite{neo4jlabs-text2cypher-2024, longwell-etal-2024-triple-augmented, yang-etal-2026-enhancing, cypherbench-megagonlabs-2024, dabramo-etal-2025-investigating}.

\paragraph{Schema Retrieval Approaches.}
Regarding the inclusion of schema information in prompts, most existing works \cite{neo4jlabs-text2cypher-2024, cypherbench-megagonlabs-2024, longwell-etal-2024-triple-augmented, dabramo-etal-2025-investigating} either assume that relevant schema components are already known or simply provide the entire database schema without any filtering or pruning. This approach, while straightforward, can lead to excessively large prompts that may exceed the model's maximum context length or introduce irrelevant information that degrades generation quality.

To address problems \ref{itm:identify}, \ref{itm:discovery}, and \ref{itm:linking} in Section \ref{sec:introduction}, only a few works have proposed techniques for \textbf{targeted schema retrieval}. Some approaches \cite{yang-etal-2026-enhancing, grasp-walter-bast-2025, grasp-entity-linking-2025} combine lexical matching methods (e.g., Levenshtein distance, fuzzy string matching, token-based or prefix matching) for schema element identification with semantic similarity techniques (e.g., BM25 sparse retrieval or dense embeddings from Sentence Transformers \cite{reimers-2019-sentence-bert}) for instance retrieval.

\paragraph{Trade-offs in Similarity Techniques.}
These retrieval methods, however, present inherent trade-offs. Lexical matching excels at identifying exact matches or minor spelling variations in short terms but becomes less effective for longer phrases where semantic understanding is crucial (e.g., handling synonymy, polysemy, and other lexical-semantic relationships). Conversely, contextual dense embeddings from transformer models excel at capturing semantic similarity and scale efficiently to large databases through vector stores. However, they can be inefficient or ineffective for short, information-dense terms such as node labels (\verb|DataCenter|, \verb|PointOfInterest|), relationship types (\verb|PARENT_OF|, \verb|DEPENDS_ON|), or property keys (\verb|created_at|, \verb|description|) which lack the rich explicit context needed to generate discriminative contextual representations useful for retrieval purposes.

Static Word Embedding (SWEs) models like FastText \cite{bojanowski-etal-2017-enriching} address this limitation by providing context-free, semantically rich representations that handle out-of-vocabulary terms through subword information. However, their prohibitive memory requirements (vocabularies with millions of tokens) make them impractical for large-scale deployment.

\textbf{Model2Vec} \cite{tulkens-2024-model2vec} offers an optimal solution for this use case. It produces static word embeddings by distilling knowledge from large Sentence Transformer models through output embedding aggregation and dimensionality reduction via PCA \cite{shlens-2014-tutorial-pca}, eliminating the need for the original model's heavy neural architecture. This distillation process yields compact, semantically rich representations that retain the semantic quality of their teacher models while requiring only vocabulary-level storage and enabling inference through simple vocabulary lookups and mean pooling. For our scenario, Model2Vec provides an ideal balance: it captures rich semantic information like contextual models, operates efficiently on short terms like static embeddings, yet avoids the computational overhead of transformer inference and the memory footprint of traditional SWEs.

\paragraph{Algorithmic vs. Agentic Architectures.}
Text-to-query solutions can be further categorized based on their architectural paradigm. \textbf{Algorithmic approaches}, such as Yang et al. \cite{yang-etal-2026-enhancing}, follow predefined pipelines: first performing entity linking via lexical matching at the schema level and hybrid lexical-semantic matching at the instance level, then extracting relevant schema portions for query construction. While deterministic, these approaches face several limitations: (i) they rely on strong, domain-specific assumptions that may not generalize, (ii) they lack robust disambiguation mechanisms when multiple candidates match the same concept, and (iii) they treat relationship linking as secondary to entity linking, imposing rigid dependencies that may not hold in real-world scenarios.

In contrast, \textbf{agentic approaches} \cite{yao-etal-2023-react, schick-etal-2023-toolformer, patil-etal-2023-gorilla}, such as the framework proposed by Walter et al. \cite{grasp-walter-bast-2025, grasp-entity-linking-2025}, equip LLM agents with a set of tools to iteratively address all sub-tasks \ref{itm:discovery}, \ref{itm:linking}, and \ref{itm:representing} from Section \ref{sec:introduction}. These tools include pre-built indexes for lexical and semantic search, as well as direct database query capabilities for grounding and query formulation. This architecture enables dynamic tool selection based on context, allowing the agent to adapt its retrieval strategy to specific database characteristics and handle ambiguities more flexibly than rigid algorithmic pipelines. However, ensuring efficiency and effectiveness at scale, particularly for large databases or complex queries, remains an open challenge requiring careful system design and optimization.
